% TEMPLATE-MILC-v3.tex - plantilla maestra para todas las guias MILC v3
% Ver scripts/generadores/build-guia-milc.py para la lista completa de
% claves esperadas. El script valida que no queden placeholders sin
% reemplazar antes de escribir el archivo final.

\documentclass[11pt,letterpaper]{article}
\usepackage[margin=1.75cm]{geometry}
\usepackage[table]{xcolor}
\usepackage{fontspec}
\usepackage[spanish]{babel}
\usepackage{tikz}
\usepackage[most]{tcolorbox}
\usepackage{tabularx}
\usepackage{array}
\usepackage{fancyhdr}
\usepackage{titlesec}
\usepackage{ragged2e}
\usepackage{enumitem}
\usepackage{microtype}

\setmainfont{Helvetica Neue}
\setsansfont{Helvetica Neue}
\setlength{\parindent}{0pt}
\setlength{\parskip}{6pt}
\hyphenpenalty=200
\exhyphenpenalty=200
\pretolerance=400
\tolerance=2000
\setlength{\emergencystretch}{1em}
\renewcommand{\arraystretch}{1.25}
\setlist{leftmargin=5mm,itemsep=1mm,topsep=1mm}
\newcolumntype{Y}{>{\RaggedRight\arraybackslash}X}

\definecolor{milcMagenta}{HTML}{E6007E}
\definecolor{milcVino}{HTML}{5A0038}
\definecolor{milcTurquesa}{HTML}{0093A5}
\definecolor{milcVerde}{HTML}{58A923}
\definecolor{milcMostaza}{HTML}{E5B400}
\definecolor{milcOcre}{HTML}{B86600}
\definecolor{milcNegro}{HTML}{241816}
\definecolor{milcGris}{HTML}{F7F7F4}
\definecolor{milcLinea}{HTML}{E8E2DD}
\definecolor{milcGradePrimary}{HTML}{84CC16}
\definecolor{milcGradeDark}{HTML}{4D7C0F}
\definecolor{milcGradeSoft}{HTML}{ECFCCB}
\definecolor{milcGradeAccent}{HTML}{CFE7FF}
\definecolor{milcGradeText}{HTML}{FFFFFF}
\color{milcNegro}

\pagestyle{fancy}
\fancyhf{}
\lhead{\small Tecnología e Informática · Grado 7}
\rhead{\small Guía 14 · MILC v3}
\cfoot{\small\thepage}
\renewcommand{\headrulewidth}{0pt}

\newtcolorbox{titlebox}[1]{
  enhanced,colback=#1,colframe=#1,arc=7mm,boxrule=0pt,
  left=7mm,right=7mm,top=3mm,bottom=3mm,
  before skip=8pt,after skip=8pt,
  fontupper=\sffamily\bfseries\Large\color{white}
}

\newtcolorbox{softbox}[3]{
  enhanced,colback=#2,colframe=#1,borderline west={4pt}{0pt}{#1},
  arc=2mm,boxrule=.4pt,left=4mm,right=4mm,top=3mm,bottom=3mm,
  before skip=6pt,after skip=8pt,title={#3},coltitle=white,
  colbacktitle=#1,fonttitle=\sffamily\bfseries,fontupper=\RaggedRight,
  attach boxed title to top left={xshift=3mm,yshift=-2mm},
  boxed title style={arc=2mm, boxrule=0pt}
}

\newtcolorbox{drawbox}[2]{
  enhanced,colback=white,colframe=#1,arc=4mm,boxrule=1pt,
  left=5mm,right=5mm,top=4mm,bottom=4mm,before skip=8pt,after skip=8pt,
  title={#2},coltitle=white,colbacktitle=#1,fonttitle=\sffamily\bfseries,
  fontupper=\RaggedRight,
  attach boxed title to top left={xshift=4mm,yshift=-2mm},
  boxed title style={arc=3mm, boxrule=0pt}
}

\newtcolorbox{infoband}[1]{
  enhanced,colback=#1!8,colframe=#1!50,arc=2mm,boxrule=.5pt,
  left=4mm,right=4mm,top=2mm,bottom=2mm,before skip=4pt,after skip=4pt,
  fontupper=\small\RaggedRight
}

% Puente narrativo entre secciones (contrato editorial Conéctate).
% Da continuidad: "Acabas de X. Ahora vas a Y porque Z."
% Visualmente: banda izquierda en color del grado, texto en itálica pequeña.
\newtcolorbox{puentebox}{
  enhanced,colback=milcGradeSoft,colframe=milcGradeDark,
  borderline west={3pt}{0pt}{milcGradeDark},
  arc=0mm,boxrule=0pt,left=5mm,right=4mm,top=2.5mm,bottom=2.5mm,
  before skip=4pt,after skip=8pt,
  fontupper=\itshape\small\color{milcNegro}\RaggedRight
}
\newcommand{\puente}[1]{%
  \begin{puentebox}%
    \textcolor{milcGradeDark}{\textbf{\textrightarrow}}\hspace{2mm}#1%
  \end{puentebox}%
}

\newcommand{\linead}[1]{\noindent\color{milcLinea}\rule{#1}{0.5pt}\color{milcNegro}}
\newcommand{\respuesta}[1]{\par\vspace{2mm}\foreach \n in {1,...,#1}{\noindent\color{milcLinea}\rule{\linewidth}{0.5pt}\par\vspace{5mm}}\color{milcNegro}}
\newcommand{\checkbox}{\textcolor{milcGradeDark}{\fbox{\phantom{x}}}\hspace{5pt}}

\newcommand{\phasecard}[4]{
\begin{tcolorbox}[enhanced,colback=#3,colframe=#2,arc=3mm,boxrule=.5pt,height=3.05cm,valign=center,halign=center,left=1.5mm,right=1.5mm,top=2mm,bottom=2mm]
{\sffamily\bfseries\fontsize{22}{24}\selectfont\textcolor{#2}{#1}}\par
{\sffamily\bfseries\small #4}
\end{tcolorbox}
}

\begin{document}
\thispagestyle{empty}
\begin{tikzpicture}[remember picture,overlay]
  \fill[white] (current page.south west) rectangle (current page.north east);
  \path[fill=milcGradePrimary,rounded corners=16mm]
    ([xshift=.55cm,yshift=-.55cm]current page.north west) rectangle
    ([xshift=-.55cm,yshift=3.65cm]current page.south east);

  \node[anchor=north west,text=milcGradeText] at ([xshift=2.0cm,yshift=-1.85cm]current page.north west)
    {\sffamily\bfseries\fontsize{14}{16}\selectfont GRADO};
  \node[anchor=north west,text=milcGradeText,opacity=.82] at ([xshift=2.0cm,yshift=-2.75cm]current page.north west)
    {\sffamily\bfseries\fontsize{9}{11}\selectfont SERIE GUÍAS · TECNOLOGÍA E INFORMÁTICA};

  \begin{scope}[shift={([xshift=-3.05cm,yshift=-2.55cm]current page.north east)}]
    \fill[white,opacity=.80] (-.62,-.52) rectangle (.62,.52);
    \draw[milcGradeText,line width=1pt] (-.62,-.52) rectangle (.62,.52);
    \fill[milcGradeDark] (-.42,-.38) rectangle (-.22,.06);
    \fill[milcGradeAccent] (-.10,-.38) rectangle (.10,.24);
    \fill[milcGradePrimary!60!black] (.22,-.38) rectangle (.42,.38);
  \end{scope}

  \node[anchor=west,text=milcGradeText] at ([xshift=1.9cm,yshift=-12.85cm]current page.north west)
    {\sffamily\bfseries\fontsize{124}{124}\selectfont 7};
  \draw[milcGradeText,line width=1.7pt]
    ([xshift=4.52cm,yshift=-11.68cm]current page.north west) circle (.13cm);

  \node[anchor=west,text=milcGradeText,text width=8.7cm,align=left] at ([xshift=10.35cm,yshift=-11.6cm]current page.north west)
    {
     {\sffamily\bfseries\fontsize{19}{22}\selectfont Diagramas\\de flujo ---\\dibujar\\el algoritmo}\\[4mm]
     {\sffamily\bfseries\fontsize{23}{26}\selectfont Guía 14-7-TIC}\\[2mm]
     {\sffamily\fontsize{13}{16}\selectfont Período 2 · Algoritmia y pensamiento computacional}\\[5mm]
     {\sffamily\bfseries\fontsize{11}{13}\selectfont Institución Educativa Sor María Juliana}\\[1mm]
     {\sffamily\fontsize{10}{12}\selectfont Autor: Dr. Álvaro Cárdenas Orozco}
    };

  \path[fill=milcGradeSoft,rounded corners=7mm]
    ([xshift=1.35cm,yshift=.85cm]current page.south west) rectangle
    ([xshift=-1.35cm,yshift=4.25cm]current page.south east);
  \draw[milcGradeDark,line width=.65pt,rounded corners=7mm]
    ([xshift=1.35cm,yshift=.85cm]current page.south west) rectangle
    ([xshift=-1.35cm,yshift=4.25cm]current page.south east);
  \node[anchor=center,text=milcNegro,text width=17.0cm,align=center] at ([yshift=2.55cm]current page.south)
    {\sffamily\fontsize{10}{12}\selectfont Metodología MILC · Apertura ancestral · Pensamiento computacional · Triángulo Dussel-Estoicismo-Floridi\\
    \textcolor{milcGradeDark}{\bfseries Producto:} Un \textbf{diagrama de flujo dibujado en el cuaderno} de un algoritmo de tu vida (puede ser el de la S2 o uno nuevo). El diagrama incluye: \textbf{(a) 5 símbolos estándar bien dibujados} (óvalo de inicio/fin, rectángulo de proceso, rombo de decisión, paralelogramo de entrada/salida, flechas direccionales); \textbf{(b) mínimo 1 decisión (rombo con SÍ/NO)}; \textbf{(c) tabla de los 5 símbolos} con su significado al lado del diagrama; \textbf{(d) firma}.};
\end{tikzpicture}
\mbox{}
\newpage

\section*{Apertura: Diagramas de flujo --- dibujar el algoritmo}

\begin{softbox}{milcGradeDark}{milcGradeSoft}{Saber ancestral}
\textbf{Los albañiles de Cartago, antes de levantar una pared, hacían un dibujo a mano en una hoja: el plano.} El plano no era la pared, era \emph{una representación visual} de cómo se levantaría: el ancho, la altura, dónde van las ventanas, dónde la puerta. \textbf{Sin plano, los albañiles no sabían por dónde empezar; con plano, todos sabían qué seguía}. El plano \emph{traducía la idea} (que estaba en la cabeza del arquitecto) \emph{en algo que cualquier albañil podía ejecutar}. En la cocina pasaba parecido: cuando una abuela tenía una receta complicada (como un pastel de cumpleaños), \textbf{hacía un esquema en papel} antes de empezar: \emph{``primero esto, después esto, si el horno está caliente entonces esto otro, si no, esperar 10 minutos''}. Ese esquema con flechas era un \textbf{diagrama de flujo} avant la lettre. Hoy aprendes los símbolos estándar para hacer diagramas de flujo de tus algoritmos.
\end{softbox}

\begin{softbox}{milcTurquesa}{milcGris}{Saber contemporáneo}
Un \textbf{diagrama de flujo} (en inglés \emph{flowchart}) es la representación visual de un algoritmo. Usa \textbf{símbolos geométricos estándar} (cada forma significa algo específico) y \textbf{flechas direccionales} para mostrar el orden y la dirección del proceso. Es como el plano del albañil pero para algoritmos. \textbf{Cinco símbolos básicos cubren el 90\% de los diagramas:} \emph{óvalo} (inicio/fin), \emph{rectángulo} (acción o proceso), \emph{rombo} (decisión, con 2 caminos SÍ/NO), \emph{paralelogramo} (entrada o salida de datos), \emph{flecha} (dirección del flujo). \textbf{Por qué se usan:} (a) \emph{visualizar} es más rápido que leer pasos; (b) \emph{detectar errores} se hace más fácil viendo el dibujo; (c) \emph{comunicar} a otra persona el algoritmo sin necesidad de explicarlo de palabra. Los diagramas de flujo se usan en programación, ingeniería, medicina, gestión de negocios y muchas otras profesiones.
\end{softbox}

\begin{softbox}{milcMostaza}{milcGris}{Pregunta-puente}
\textit{Si tuvieras que explicar a tu mamá \emph{cómo decidir qué ponerse de ropa cada mañana} (con todas sus decisiones: ¿llueve o no? ¿hace frío o calor? ¿es para colegio o paseo?), ¿\textbf{te saldría mejor en lista de pasos o en dibujo con flechas}?}
\end{softbox}

\begin{softbox}{milcVerde}{milcGris}{Saber y saber hacer}
Al terminar la clase: (1) podrás \textbf{identificar} los 5 símbolos estándar de un diagrama de flujo; (2) sabrás \textbf{aplicar} cada símbolo a su función; (3) podrás \textbf{evaluar} si un diagrama está bien dibujado; (4) habrás \textbf{creado} tu primer diagrama de flujo de un algoritmo real.
\end{softbox}

\small
\begin{tabularx}{\textwidth}{>{\bfseries}p{3.0cm}Y}
\rowcolor{milcGradePrimary!12}Autor & Dr. Álvaro Cárdenas Orozco\\
DBA & Construye algoritmos y diagramas de flujo para resolver problemas paso a paso (MEN, T\&I 7°)\\
Referentes & Saber ancestral del tejedor · Pensamiento computacional · Scratch\\
Producto final & Un \textbf{diagrama de flujo dibujado en el cuaderno} de un algoritmo de tu vida (puede ser el de la S2 o uno nuevo). El diagrama incluye: \textbf{(a) 5 símbolos estándar bien dibujados} (óvalo de inicio/fin, rectángulo de proceso, rombo de decisión, paralelogramo de entrada/salida, flechas direccionales); \textbf{(b) mínimo 1 decisión (rombo con SÍ/NO)}; \textbf{(c) tabla de los 5 símbolos} con su significado al lado del diagrama; \textbf{(d) firma}.\\
\end{tabularx}
\normalsize

\puente{Hoy haces 4 cosas: identificas los 5 símbolos, los aplicas a un mini-ejemplo, dibujas tu propio diagrama, lo verificas.}

\begin{titlebox}{milcGradeDark}
Ruta MILC: así vamos a aprender
\end{titlebox}
\begin{tabularx}{\textwidth}{>{\centering\arraybackslash}X>{\centering\arraybackslash}X>{\centering\arraybackslash}X>{\centering\arraybackslash}X}
\phasecard{1}{milcVerde}{milcGris}{Escuta\\[-1mm]\small Observo, escucho y pregunto.} &
\phasecard{2}{milcTurquesa}{milcGris}{Sistematización\\[-1mm]\small Pensamiento computacional.} &
\phasecard{3}{milcMagenta}{milcGris}{Praxis\\[-1mm]\small Construyo y aplico.} &
\phasecard{4}{milcVino}{milcGris}{Evaluación\\[-1mm]\small Reflexión liberadora.}
\end{tabularx}


\puente{Empezamos viendo un diagrama de flujo real para reconocer sus partes.}

\begin{titlebox}{milcVerde}
Fase 1 · Escuta
\end{titlebox}

\begin{softbox}{milcVerde}{milcGris}{Escena para escuchar}
\textbf{Actividad 1 · IDENTIFICA --- Lee un diagrama de flujo} (10 min · individual). Imagina o el profe te muestra el siguiente diagrama de flujo simple del algoritmo \emph{``¿voy al colegio?''}: \emph{(óvalo) INICIO → (rombo) ¿Es día de semana? SÍ → (rombo) ¿Es hora de salir? SÍ → (rectángulo) Vestirme con uniforme → (rectángulo) Desayunar → (rectángulo) Tomar el bus → (óvalo) FIN. Si ¿Es día de semana? = NO → (rectángulo) Dormir más → (óvalo) FIN. Si ¿Es hora de salir? = NO → (rectángulo) Esperar → vuelve al rombo}. \textbf{En el cuaderno responde:} \emph{(1) ¿Qué representan los óvalos?}. \emph{(2) ¿Qué representan los rombos?}. \emph{(3) ¿Qué representan los rectángulos?}. \emph{(4) ¿Cuántos caminos finales posibles hay en este diagrama?}. \emph{(5) ¿Las flechas tienen una sola dirección o varias?}.
\end{softbox}

\checkbox Imaginé o leí el diagrama de flujo.\\
\checkbox Respondí las 5 preguntas con mi intuición.\\
\checkbox Pensé en los significados de cada símbolo.

\begin{infoband}{milcVerde}
\textbf{Lo que va al cuaderno (Actividad 1 · IDENTIFICA).} Título: «Actividad 1 --- Lectura de un diagrama». \textbf{Formato:} 5 respuestas cortas. \textbf{Extensión:} 5 líneas. \textbf{Sabes que terminaste cuando} reconoces que cada forma tiene un significado específico.
\end{infoband}

\puente{Después aprendes los 5 símbolos estándar + las 4 reglas para dibujar bien.}

\begin{titlebox}{milcTurquesa}
Fase 2 · Sistematización con pensamiento computacional
\end{titlebox}

\begin{softbox}{milcTurquesa}{milcGris}{Cuatro pilares aplicados al tema}
Tu intuición acertó: cada símbolo significa algo. Ahora aprendes los 5 símbolos estándar + las 4 reglas para dibujar diagramas claros.
\end{softbox}

\small
\begin{tabularx}{\textwidth}{>{\bfseries\RaggedRight\arraybackslash}p{3.6cm}Y}
\rowcolor{milcTurquesa!90}\color{white}Pilar & \color{white}\bfseries Aplicación al tema\\
1. Descomposición & \textbf{Los 5 símbolos estándar.} \textbf{(1) Óvalo (o pastilla)}: \emph{Inicio} y \emph{Fin} del algoritmo. Va al comienzo y al final. \textbf{(2) Rectángulo}: \emph{Proceso} o \emph{acción}. Una instrucción que se ejecuta. Adentro va un verbo (\emph{``Mojar manos''}, \emph{``Multiplicar por 2''}). \textbf{(3) Rombo}: \emph{Decisión} (\emph{``¿es par?''}, \emph{``¿está lloviendo?''}). De un rombo salen 2 caminos: uno por SÍ y otro por NO. \textbf{(4) Paralelogramo}: \emph{Entrada} (datos que recibe el algoritmo) o \emph{Salida} (resultado que muestra). En lenguaje natural: \emph{``Ingresar nombre''}, \emph{``Mostrar resultado''}. \textbf{(5) Flecha}: dirección del flujo. Conecta los símbolos en el orden de ejecución.\\
2. Reconocimiento de patrones & \textbf{Las 4 reglas para dibujar bien.} (1) \textbf{Un solo INICIO y un solo FIN.} El algoritmo arranca en un lugar y termina en un lugar. Si tiene varios fines, cada uno se simboliza pero la lógica debe ser clara. (2) \textbf{Flechas en una sola dirección.} Las flechas van de un símbolo a otro siguiendo la lógica. NO se dibujan flechas que vayan hacia atrás sin razón. (3) \textbf{Cada rombo tiene exactamente 2 salidas: SÍ y NO}. Etiqueta cada flecha que sale del rombo. (4) \textbf{Símbolos del mismo tamaño aproximado.} Para que se vea profesional. Si haces algunos gigantes y otros chiquitos, se ve desorganizado.\\
3. Abstracción & \textbf{Cómo se lee un diagrama (paso a paso).} Empiezas en el óvalo INICIO. Sigues la flecha al siguiente símbolo. Si es rectángulo, ejecutas la acción mentalmente. Si es paralelogramo, ingresas/obtienes el dato. Si es rombo, evalúas la pregunta: si la respuesta es SÍ, vas por la flecha del SÍ; si es NO, por la del NO. Continúas hasta llegar al óvalo FIN. \textbf{Importante:} en un diagrama bien hecho, NO te puedes perder. Siempre hay un solo camino que seguir desde cada símbolo.\\
4. Algoritmo conceptual & \textbf{Software opcional para hacer diagramas digitales.} Para 7° grado, dibujar a mano en el cuaderno es suficiente. Pero si quieres en digital: \emph{draw.io} (\texttt{app.diagrams.net}, gratis, en navegador), \emph{Lucidchart} (versión gratis básica), \emph{Microsoft Visio} (parte de Office 365). Lo importante es saber los símbolos; las herramientas varían.\\
\end{tabularx}
\normalsize

\begin{drawbox}{milcTurquesa}{5 símbolos + 4 reglas}
\textbf{5 símbolos:} \\[1mm]
\textbf{1.} Óvalo --- Inicio / Fin. \\[1mm]
\textbf{2.} Rectángulo --- Proceso. \\[1mm]
\textbf{3.} Rombo --- Decisión (SÍ/NO). \\[1mm]
\textbf{4.} Paralelogramo --- Entrada / Salida. \\[1mm]
\textbf{5.} Flecha --- Dirección. \\[1mm]
\textbf{4 reglas:} 1 inicio y 1 fin · flechas con dirección · rombo con 2 salidas etiquetadas · símbolos del mismo tamaño.
\end{drawbox}

\begin{softbox}{milcMostaza}{milcGris}{Errores comunes y su corrección}
\textbf{Errores frecuentes al dibujar diagramas.} (1) \emph{Olvidar el óvalo de INICIO o FIN} --- el diagrama queda incompleto. (2) \emph{Rombo sin etiquetar SÍ/NO} --- no se sabe por dónde va cada flecha. (3) \emph{Flechas que se cruzan sin razón} --- queda enredado, difícil de leer. Trata de dibujar de arriba a abajo o de izquierda a derecha. (4) \emph{Confundir símbolos} --- proceso en rombo, decisión en rectángulo. Cada forma tiene su uso. (5) \emph{Diagrama enorme con 30 símbolos} --- si tu algoritmo es complejo, mejor dividirlo en sub-diagramas más simples (técnica de descomposición de la S3).
\end{softbox}

\begin{infoband}{milcTurquesa}
\textbf{Lo que va al cuaderno (Actividad 2 · EXPLICA).} Título: «Actividad 2 --- 5 símbolos + 4 reglas». \textbf{Formato:} bloque A con tabla de 5 filas, 2 columnas (símbolo dibujado + significado). Bloque B con las 4 reglas. \textbf{Extensión:} 5 filas + 4 reglas. \textbf{Sabes que terminaste cuando} puedes dibujar los 5 símbolos de memoria con su significado.
\end{infoband}

\puente{Con todo claro, dibujas tu diagrama de flujo de un algoritmo de tu vida.}

\begin{titlebox}{milcMagenta}
Fase 3 · Praxis con pensamiento computacional
\end{titlebox}

\begin{softbox}{milcMagenta}{milcGris}{Construyendo el producto con los cuatro pilares}
\textbf{Actividad 3 · CREA --- Tu diagrama de flujo} (28 min · individual). Vas a dibujar un diagrama de flujo de un algoritmo. Puede ser el de la S2 (tu primer algoritmo) o uno nuevo más interesante. Sugerencia: hazlo de algo con \emph{al menos 1 decisión} (por ejemplo \emph{``decidir qué hacer si llueve o no llueve''} cuando vas al colegio).
\end{softbox}

\small
\begin{tabularx}{\textwidth}{>{\bfseries\RaggedRight\arraybackslash}p{3.6cm}Y}
\rowcolor{milcMagenta!90}\color{white}Pilar & \color{white}\bfseries Aplicación al producto\\
Descomposición & \textbf{Paso 1 --- Escoge el algoritmo a diagramar.} Si quieres, toma el de la S2. Si no, escoge uno nuevo con \emph{al menos una decisión}. Sugerencias: \emph{``Decidir qué ropa ponerme hoy''}, \emph{``Resolver una pelea con un hermano''}, \emph{``Encontrar mi celular cuando se pierde''}, \emph{``Decidir si veo TikTok o estudio''}. Si tiene 2 decisiones, mejor.\\
Patrones & \textbf{Paso 2 --- Lista de pasos en lenguaje natural.} Antes de dibujar, escribe los pasos del algoritmo en una lista numerada. Identifica cuáles son acciones (rectángulo), cuáles son decisiones (rombo), cuáles son entradas o salidas (paralelogramo). \emph{Esta lista es tu borrador previo}: te ayuda a no perderte al dibujar.\\
Abstracción & \textbf{Paso 3 --- Dibuja el diagrama.} En una página del cuaderno (idealmente apaisada para tener más espacio), empieza con el \textbf{óvalo INICIO} en la parte superior. Después dibuja cada símbolo correspondiente, en orden, conectado con flechas. Para los rombos, etiqueta cada flecha de salida con \emph{SÍ} o \emph{NO}. Termina con el \textbf{óvalo FIN}. Usa regla si quieres para que las líneas queden derechas.\\
Algoritmo de armado & \textbf{Paso 4 --- Tabla de símbolos al lado.} En un costado de la hoja, dibuja una mini-tabla con los 5 símbolos que usaste y al lado su significado: \emph{``Óvalo = inicio o fin''}, \emph{``Rectángulo = proceso''}, etc. Esto entrena a quien lo lea y a ti mismo a no olvidarlos.\\
\end{tabularx}
\normalsize

\begin{drawbox}{milcMagenta}{Antes de cerrar la S4}
\checkbox Tengo lista de pasos antes de dibujar (borrador). \\[1mm]
\checkbox El diagrama tiene óvalo INICIO y óvalo FIN. \\[1mm]
\checkbox Hay al menos 1 rombo de decisión con SÍ/NO etiquetados. \\[1mm]
\checkbox Las flechas tienen una dirección clara. \\[1mm]
\checkbox La tabla de los 5 símbolos está al lado del diagrama. \\[1mm]
\checkbox Un compañero verificó que entiende sin explicación.
\end{drawbox}

\begin{softbox}{milcMagenta}{milcGris}{Plantilla de trabajo}
\textbf{Estructura.} \textit{Borrador previo:} lista de pasos numerada con tipo de símbolo. \textit{Diagrama:} óvalo INICIO → símbolos conectados con flechas → al menos 1 rombo SÍ/NO → óvalo FIN. \textit{Tabla lateral:} 5 símbolos con significado. \textit{Verificación:} compañero entiende. \textit{Final:} reflexión + firma.
\end{softbox}

\begin{infoband}{milcMagenta}
\textbf{Lo que va al cuaderno (Actividad 3 · CREA).} Título: «Actividad 3 --- Mi primer diagrama de flujo». \textbf{Formato:} borrador + diagrama dibujado + tabla de símbolos + reflexión. \textbf{Extensión:} 1 hoja entera. \textbf{Sabes que terminaste cuando} otro estudiante puede leer tu diagrama y entender el algoritmo sin que tú expliques.
\end{infoband}

\puente{El producto es ese diagrama + tabla de símbolos.}

\begin{titlebox}{milcMostaza}
Producto de la sesión
\end{titlebox}
\textbf{Diagrama de flujo · 5 símbolos · 1+ decisión · verificado · firmado}.
\begin{softbox}{milcMostaza}{milcGris}{Criterios de calidad}
(1) El diagrama tiene óvalo INICIO y FIN. (2) Hay al menos 1 rombo de decisión con SÍ/NO etiquetados. (3) Las flechas siguen una dirección lógica. (4) La tabla de símbolos está al lado. (5) El diagrama es entendible sin explicación verbal.
\end{softbox}

\puente{Antes de salir, intercambia el diagrama con un compañero: ¿él entiende sin que tú expliques?}

\begin{titlebox}{milcVino}
Fase 4 · Evaluación liberadora — cinco dimensiones
\end{titlebox}

\begin{softbox}{milcVino}{milcGris}{Sentido de la evaluación}
Evaluar no es solo revisar una nota. Es reconocer qué aprendí, qué cambió en mí, qué emociones manejé, cómo me relaciono con la ciudad y la comunidad, y qué le dejaré a quien viene detrás.
\end{softbox}

\textbf{Mi reflexión liberadora en cinco dimensiones.} Completa cada frase iniciada:

\small
\begin{tabularx}{\textwidth}{>{\bfseries\RaggedRight\arraybackslash}p{3.4cm}Y>{\RaggedRight\arraybackslash}p{4.6cm}}
\rowcolor{milcVino}\color{white}Dimensión & \color{white}\bfseries Mi respuesta (frame starter) & \color{white}\bfseries Algo que descubrí\\
1. Personal & Lo que más cambió en mí fue \linead{6.5cm} & \linead{4.2cm}\\
2. Emocional & La emoción más fuerte fue \linead{2.5cm} y la regulé \linead{3cm} & \linead{4.2cm}\\
3. Ciudadana & En democracia esto importa porque \linead{6cm} & \linead{4.2cm}\\
4. Local (Cartago) & En mi barrio sería útil para \linead{6.5cm} & \linead{4.2cm}\\
5. Intergeneracional & A mi abuela/abuelo le diría \linead{6.5cm} & \linead{4.2cm}\\
\end{tabularx}
\normalsize

\begin{infoband}{milcVino}
\textbf{Para la próxima sustentación.} Vuelve a esta tabla en seis meses. Verás que las respuestas cambian — no porque lo aprendido cambie, sino porque tú cambias. La evaluación liberadora es una conversación contigo a través del tiempo.
\end{infoband}


\puente{Cerramos con 3 voces sobre por qué un dibujo claro vale más que mil palabras.}

\begin{titlebox}{milcGradeDark}
Triángulo de pensamiento · cierre filosófico
\end{titlebox}

\begin{softbox}{milcVino}{milcGris}{Dussel · lente del nosotros}
\textit{````Los símbolos compartidos son lenguaje universal. El que los conoce habla con cualquier persona del mundo, en cualquier época.''''}

\textbf{Los símbolos del diagrama de flujo (óvalo, rectángulo, rombo) son universales}. Un ingeniero en Tokio, una médica en Cartago y un programador en Berlín los entienden igual. Cuando aprendes estos símbolos, te conviertes en miembro de una \emph{comunidad global de pensadores estructurados}. La abuela albañil que dibujaba el plano de la pared participaba de esa misma tradición: \emph{trazar lo abstracto en formas concretas que cualquiera entiende}.

\textbf{¿Qué otros lenguajes simbólicos universales conozco (matemáticas, música, señales de tráfico)? ¿Qué tienen en común?}
\end{softbox}
\linead{\linewidth}\\[1mm]
\linead{\linewidth}

\begin{softbox}{milcMostaza}{milcGris}{Estoicismo · Marco Aurelio (emperador que usaba diagramas para planear campañas) · cuidado interior}
\textit{````Un dibujo claro vale más que mil palabras vagas. Quien dibuja antes de hablar, ahorra confusiones.''''}

Marco Aurelio planeaba campañas militares con \textbf{mapas y esquemas}, no solo con explicaciones verbales. Sabía que un dibujo \emph{transmite información compleja en segundos} mientras que una explicación verbal tarda minutos y abre lugar a malentendidos. \emph{Cuando explicas tu plan a otros, primero dibuja, después habla}. Tu equipo te lo agradece.

\textbf{¿Cuántas veces he intentado explicar algo complicado solo con palabras cuando un dibujo lo haría rápido?}
\end{softbox}
\linead{\linewidth}\\[1mm]
\linead{\linewidth}

\begin{softbox}{milcTurquesa}{milcGris}{Floridi · lente de la infoesfera}
\textit{````El siglo XXI privilegia el pensamiento visual. Quien sabe diagramar, comunica más rápido que quien solo sabe redactar.''''}

Las redes profesionales modernas (LinkedIn, presentaciones, propuestas de negocio) valoran cada vez más \textbf{el pensamiento visual}. \emph{Diagramas, infografías, mapas mentales, gráficos}: el que sabe representarlos comunica más eficazmente. Aprender diagramas de flujo en 7° grado es entrenamiento para esa habilidad adulta. No es solo para programadores: es para cualquier profesional moderno.

\textbf{¿En qué situaciones de mi vida un diagrama claro me ahorraría malentendidos?}
\end{softbox}
\linead{\linewidth}\\[1mm]
\linead{\linewidth}

\begin{titlebox}{milcVino}
Compromiso final
\end{titlebox}

\small
\begin{tabularx}{\textwidth}{>{\RaggedRight\arraybackslash}p{3.0cm}YYY}
\rowcolor{milcVino}\color{white}\bfseries Criterio & \color{white}\bfseries Lo logré & \color{white}\bfseries En proceso & \color{white}\bfseries Necesito apoyo\\
Comprensión del tema & Explico ideas centrales con mis palabras. & Reconozco ideas pero falta precisión. & Necesito repasar conceptos base.\\
Pensamiento computacional & Apliqué los 4 pilares al diseñar mi producto. & Apliqué algunos pilares. & Necesito releer la fase 2.\\
Aplicación y producto & Mi producto cumple los criterios y se puede revisar. & Producto funcional con ajustes pendientes. & Aún en construcción.\\
Reflexión 5 dimensiones & Respondí cada dimensión con honestidad. & Respondí algunas con detalle. & Mis respuestas son muy generales.\\
Triángulo de pensamiento & Conecté las 3 voces con mi trabajo real. & Conecté al menos una. & No relaciono las voces con mi proceso.\\
\end{tabularx}
\normalsize

\textbf{Hoy entendí que:}\\[1mm]
\linead{\linewidth}\\[1mm]
\linead{\linewidth}

\textbf{Mi compromiso para la próxima sesión es:}\\[1mm]
\linead{\linewidth}\\[1mm]
\linead{\linewidth}

\begin{softbox}{milcMostaza}{milcGris}{Cita de despedida · Marco Aurelio}
\textit{``El obstáculo en el camino se vuelve el camino. Lo que estorba al camino se vuelve camino.''}\\[1mm]
Cada error de hoy, bien observado, se convierte en pieza del camino que pisarás mañana.
\end{softbox}

\small
\begin{tabularx}{\textwidth}{>{\bfseries}p{3.6cm}Y}
\rowcolor{milcGradeDark!12}Nombre completo: & \linead{10cm}\\
Fecha: & \linead{4cm} \quad Curso/grupo: \linead{4cm}\\
Mi firma: & \linead{9cm}\\
\end{tabularx}
\normalsize

\begin{infoband}{milcGradeDark}
\textbf{Para la próxima sesión.} Esta semana, dibuja al menos 1 diagrama de flujo más (puede ser de algo de otra materia o de tu rutina diaria). La próxima clase aprendes \textbf{variables}: las cajitas con etiqueta que guardan información en los algoritmos. Es el primer componente fundamental de toda programación.
\end{infoband}

\end{document}
